\section{Marked sums and analytic inputs}\label{sec:framework}

Put $L=\log x$ and $\e(t)=\exp(2\pi it)$. A constant denoted by
$C$ may change between estimates. Every exponent of $L$ is fixed before
$x$ tends to infinity. The notation $O_A(L^{-A})$ means that each fixed
accuracy $A>0$ is available, after the stated choices of parameters.

\subsection{Long slots, inert bands, and active groups}

Fix $d\geq1$ and $\eps>0$. Let $I_1,\ldots,I_d$ be intervals of
primes whose lower endpoints are at least $x^\eps$ and whose upper
endpoints have product at most $x^{1-\eps}$. Tuples are ordered and
their numerical prime entries are distinct. Write
\begin{equation}\label{eq:tuple}
 T_x(n)=\sum_{\boldsymbol\ell}\ind_{\ell\mid n},\qquad
 \ell=\ell_1\cdots\ell_d,\qquad C_x=\sum_{\boldsymbol\ell}\frac1\ell.
\end{equation}
On each fixed range $n\leq x^C$, $T_x(n)=O_{C,d,\eps}(1)$.
Also $C_x=O_{d,\eps}(1)$ by the harmonic prime estimate.

An inert array consists of $K$ disjoint full prime intervals
$[\exp(L^{a_i}),\exp(2L^{a_i}))$, with fixed $.1<a_i<.2$.
For damping $0<q<1$, its weight is
\begin{equation}\label{eq:inert}
 \W(n)=\prod_{i=1}^K q^{\omega_i(n)-1}\frac{\omega_i(n)}{V_i},
 \quad \omega_i(n)=\sum_{p\in\mathcal P_i}\ind_{p\mid n},
 \quad V_i=\sum_{p\in\mathcal P_i}\frac1p.
\end{equation}
A factor is zero when $\omega_i(n)=0$. For fixed $K,q$, the weight
is bounded. Its independent divisor-model mean $m_{\W}$ is obtained
by replacing the divisibility indicators by independent Bernoulli
variables with means $1/p$. The source gives
$0<c_K\leq m_{\W}\leq C_K$ and $\E_{\rm div}\W^2\leq C_K$,
uniformly over the fixed exponents.

Independently, choose disjoint active groups $\mathcal P_g$, in two
families $\mathcal G_s,\mathcal G_b$. Each family has between fixed
positive multiples of $\log L$ groups, and
\begin{align*}
 \mathcal P_g&\subset[\exp(L^{.27}),\exp(L^{.36})],&&g\in\mathcal G_s,\\
 \mathcal P_g&\subset[\exp(L^{.39}),\exp(L^{.46})],&&g\in\mathcal G_b.
\end{align*}
Their harmonic masses lie in a fixed compact subinterval of $(0,\infty)$.
Every big group is a full prime interval. At least one small group is
also a full interval, as needed for fixed-modulus mixing. The extraction
uses full intervals for all candidate groups.

Let $n_*$ denote $n$ after removal of all powers of active primes.
For $\mathbf t=(t_g)$ and a bounded function $b$ of ordered distinct
represented divisor lists, set
\begin{equation}\label{eq:active}
 W_{\mathbf t}^{b}(n)=
 q_0^{\omega_P(n)-\sum_g t_g}\prod_g V_g^{-t_g}
 \sum_{\text{represented lists at }n} b(\mathbf p),
 \qquad 0<q_0<1.
\end{equation}
Here $\omega_P(n)=\sum_g\omega_g(n)$ counts distinct active prime
divisors, and $V_g=\sum_{p\in\mathcal P_g}1/p$.
Write $W_1$ when $t_g=1$ for all $g$ and $b=1$.
For fixed bounded list lengths and $|b|\leq1$, these weights are at
most a fixed power of $L$. The second endpoint will have the form
$G(n)=G(n_*)$ with $|G(n)|\leq L^C\tau(n_*)^C$.

\subsection{The source estimates used at the two endpoints}

We use the general physical pairing in \cite[Section 6]{D}, rather
than extending its final centered correlation theorem by substitution.
The following description fixes its measures and the parts of its
conclusion that enter the proof.

For a sufficiently large fixed integer $J$, a physical state at $n$
has $J+1$ ordered distinct prime divisors from each active group. Its
list mass is $\prod_g V_g^{-(J+1)}$, with counting measure in $n$.
In a pattern $\nu$, the source and target share $J-|C_g|$ labels in
group $g$. Small groups have $C_g=\varnothing$; big groups have
$|C_g|\leq m_0$ for a fixed $m_0$. Let $D_R$ be the shared product
and $D_C$ the product of the auxiliary free labels. Free labels have
independent laws $1/(V_gp)$. The physical shift is $D_RD_C$.
The pattern kernels depend only on free big labels and unshared big
endpoint lists. The family has total kernel mass $L^{O(1)}$, with
exponent independent of $J$.

The residual ideal-family theorem supplies a raw pattern $D_C=1$
and signed comparison patterns. Its ideal operator has arbitrarily
small logarithmic norm. The general transfer theorem then bounds the
whole physical pairing for a first vector independent of its lists at
fixed $n$, with the source supremum bound, and endpoint norms
$\leq X^{1/2}L^{C_3}$ on a position dyad $[X,2X)$.
The exponent $C_3$ must be independent of $J$.
Both endpoints retain the extra half damping
$q_0^{(\omega_P(n)-(J+1)|P|)/2}$, which is at most one on represented
states. If the first core is bounded and the second has the divisor
bound above, the required norm estimate follows from
\begin{equation}\label{eq:listnorm}
 \prod_g V_g^{-(J+1)}\sum_{\text{lists}}\frac1{\prod p}\leq1
\end{equation}
and a fixed divisor moment. Conjugate the first core in the physical
vector to cancel the conjugation in the Hilbert-space pairing.

After inserting $\Psi(n/(xD_RD_C))/n$, restoring its full Mellin
integral, and summing all position dyads, put
\begin{equation}\label{eq:strip}
 n=D_Rw,\quad n'=D_R(w+D_C),\qquad
 \frac{\Psi(n/(xD_RD_C))}{n}
   =\frac1{D_Rw}\Psi\left(\frac{w}{xD_C}\right).
\end{equation}
The remaining cutoff has no shared-label dependence. The factor
$D_R^{-1}$ makes the shared-label measures harmonic probabilities.
The source's collision and Mellin-tail errors save every fixed power
of $L$, also after bounded multipliers at the endpoints.

Every comparison has two nonempty free big products, at scales
$U_1,U_2$. Write
\begin{equation}\label{eq:scales}
 H'=U_1U_2,\quad Y=xH',\quad H=H'/L^{C_6},\quad
 T_0=LY/H=xL^{C_6+1},\qquad
 U_i\geq\exp(L^{.39}-O(1)).
\end{equation}
Expanding the source kernels reduces a comparison to bounded separate
tests on the two free tuples, with a total fixed logarithmic cost.
Its additive multiplier is small off rational arcs of denominator
$\leq L^{C_6}$ and width $H^{-1}$. The source's local Mellin estimate
reduces the first-endpoint energy on these arcs to
\begin{equation}\label{eq:energy}
 \int_{|t|\leq T_0}
 \left|Y^{-1}\sum_n \mathcal U(n)\e(hn/k)n^{it}\right|^2\dd t
 \ll_A L^{-A},\qquad k\leq L^{C_6}.
\end{equation}
Here $\mathcal U$ has a smooth cutoff at $Y$, a logarithmic-sized
twist, the proposed first core, and a factor
$\W W_{\mathbf t}^{b}$ whose test $b$ uses only big marks.
The second endpoint needs only $\sum|\mathcal V(n)|^2\ll YL^C$.
The estimates and kernel expansions in this paragraph are the input
package from \cite[Section 6]{D}; we prove \eqref{eq:energy} for the
new cores below.

The truncated divisor model for $\W W_{\mathbf t}^{b}$ has the form
\begin{equation}\label{eq:divisor}
 \sum_{a\mid n}A_a+E(n),\quad
 a\leq\exp(O(L^{.47})),\quad
 \sum_a\frac{|A_a|}{a}\leq L^C,\quad
 \sum_{n\asymp R}|E(n)|^2\ll_N RL^{-N}.
\end{equation}
It permits a zero list length when a small mark has been removed.
The constants are uniform in bounded big-mark tests.

We also use two explicit inputs from \cite[Appendix 9]{D}.
For fixed $0<\tau<\eta<1$ and $A,C>0$, there is $B_0$ such that
\begin{equation}\label{eq:longprime}
 \left|\sum_{p\in I}\chi(p)p^{-1+it}\right|\ll L^{-A}
\end{equation}
for $I\subset[N,2N]$, $x^\tau/2\leq N\leq x^\eta$,
$\operatorname{mod}(\chi)\leq L^C$, and $L^{B_0}\leq|t|\leq x^2$.
For some absolute $C_*$, the second estimate is
\begin{equation}\label{eq:logphase}
 \left|\sum_{\substack{n\in I\\n\equiv a\ (k)}}n^{it}\right|
 \ll\frac{N'}k\exp\left(-\frac{L}{(\log L)^{C_*}}\right),
\end{equation}
where $I\subset[N',2N']$, $k\leq L^C$,
$\exp(L/(\log L)^2)\leq N'\leq2x^5$, and
$\exp(L/(2(\log L)^2))\leq|t|\leq4x^3$.
The residue $a$ need not be a unit. The subinterval uniformity in both
estimates allows partial summation against the smooth amplitudes below.

\subsection{The scalar determinant and positive extraction inputs}

The coefficient criterion in \cite[Section 8]{P} applies to rough
dyadic coefficients $\alpha_m,\beta_n$, bounded by fixed powers of $L$,
on positive-power scales $H_mH_n\asymp x$. It requires
\begin{equation}\label{eq:coefficient}
 H_m^{-1}\left|\sum_m\alpha_m\chi(m)m^{it}\right|\ll_A L^{-A}
\end{equation}
for every fixed logarithmic modulus range, every fixed accuracy, and
$|t|\leq2H_mH_nL^B$. The scalar estimate in its proof bounds the
expanded Cauchy square after $mn-1=ah$, with $a$ a product of $K$
selected inert primes. An invariant bounded factor of $h$ lies outside
that square. The source prime-minus-rough coefficient meets
\eqref{eq:coefficient}. Multiplication by a character of a fixed
modulus $M$ replaces the test modulus by its least common multiple
with $M$, which is still logarithmic.

The extraction in \cite[Section 7]{D} uses a nonnegative block sieve,
prime-to-rough telescoping, and a partition into independent conditional
coin outcomes. The block sieve has relative error $O(e^{-h})$ at even
depth $h$ and uses summed progression remainders only up to
$z^{4h+2}$ when the sieve limit is $z$. For the independent divisor
model, its local densities are $1/r$ at allowed primes.
The rough proxy is
\[
 a(m)=\frac{\ind_{m>x^\gamma}\ind_{P^-(m)>W}}{V(W)\log m},
 \qquad W=\exp(L^{.24}),\quad V(W)=\prod_{p\leq W}(1-1/p).
\]
The source permits removal of all candidate prime parts from its
thresholds and logarithms at absolute cost $o(x/L)$. It also gives a
total positive proxy mass $O_{j,\gamma}(x/L)$ for each fixed number
$j$ of factors. Section~\ref{sec:extraction} states the allocation
identity and carries out the needed new recombination.
These imported estimates do not assert cancellation of $f_z(p-1)$.
